\section{Laman's theorem}
\label{sec:laman-theorem}

For $n \ge 2$ vertices, a simple graph is generically rigid in the
plane iff it contains a Laman spanning subgraph (\cref{def:isLaman}):
a $(2, 3)$-tight subgraph on the same vertex set. The two directions
are proved separately:

\begin{itemize}
  \item \emph{$\Leftarrow$ (every Laman graph is generically rigid in dimension 2).}
    A Henneberg induction on $|V|$. Base case $|V| = 2$: $K_2$ is
    generically rigid (\cref{thm:top-fin-two-isGenericallyRigidInj}).
    Inductive step: peel off a Type~I or Type~II move via
    \cref{thm:isLaman-exists-typeI-or-typeII-reverse} and apply the
    move-rigidity-preservation theorem
    (\cref{thm:typeI-isGenericallyRigidInj-two} or
    \cref{thm:typeII-isGenericallyRigidInj-two}). The induction
    maintains the strong predicate \emph{injectively} generically rigid
    (\cref{def:isGenericallyRigidInj}), then weakens at the end.
  \item \emph{$\Rightarrow$ (every generically rigid graph contains a
    Laman spanning subgraph).} The Lov\'asz--Yemini~\cite{lovaszYemini1982}
    direction, which uses the easy half of the identification of the
    rigidity matroid in dimension 2 with the $(2, 3)$-count matroid:
    row-independence of an edge set at a generic placement implies
    $(2, 3)$-sparsity of the underlying subgraph.
\end{itemize}

\subsection{The edge count in the plane}

In the plane, the edge-count bound
(\cref{thm:isGenericallyRigid-card-mul-le}) says that a generically
rigid graph on $V$ has at least $2|V| - 3$ edges, as many as a Laman
graph on $V$.

\begin{lemma}[Edge-count bound in dimension 2]
  \label{lem:isGenericallyRigid-card-mul-le-two}
  \lean{SimpleGraph.IsGenericallyRigid.card_mul_le_two}
  \leanok
  \uses{thm:isGenericallyRigid-card-mul-le}
  Let $V$ be finite, and let $G$ be generically rigid in
  dimension~$2$. Then
  \[
    2 |V| \;\le\; |E(G)| + 3.
  \]
\end{lemma}
\begin{proof}
  \leanok
  Specialize \cref{thm:isGenericallyRigid-card-mul-le} to $d = 2$,
  where $d(d+1)/2 = 3$.
\end{proof}

\subsection{The \texorpdfstring{$\Leftarrow$}{<=} direction}

The sufficiency half of Laman's theorem~\cite{laman1970} says that
every Laman graph is generically rigid in the plane. We prove it by
induction on the number of vertices, along the reverse decomposition
of \cref{sec:henneberg-decomp}. The step for each Henneberg move needs
an injective framework (\cref{sec:henneberg-rigidity}), so the
induction proves the stronger statement for injective generic
rigidity.

\begin{theorem}[Every Laman graph is generically rigid in dimension 2]
  \label{thm:isLaman-isGenericallyRigid-two}
  \lean{SimpleGraph.IsLaman.isGenericallyRigid_two}
  \leanok
  \uses{def:isLaman, def:isGenericallyRigid, def:isGenericallyRigidInj,
    thm:isLaman-exists-typeI-or-typeII-reverse,
    thm:typeI-isGenericallyRigidInj-two,
    thm:typeII-isGenericallyRigidInj-two,
    thm:top-fin-two-isGenericallyRigidInj,
    lem:isLaman-eq-top-of-card-eq-two,
    lem:isGenericallyRigidInj-iso}
  Let $V$ be finite, and let $G$ be a Laman graph on $V$. Then $G$ is
  generically rigid in dimension~$2$.
\end{theorem}
\begin{proof}
  \leanok
  We show by strong induction on $|V|$ that $G$ is injectively
  generically rigid in dimension~$2$ (\cref{def:isGenericallyRigidInj}).
  An injectively generically rigid graph is generically rigid.

  Since $|E(G)| + 3 = 2|V|$, the set $V$ has at least two elements. If
  $|V| = 2$, then $G$ is the complete graph on $V$
  (\cref{lem:isLaman-eq-top-of-card-eq-two}), which is isomorphic to
  $K_2$. The graph $K_2$ is injectively generically rigid in
  dimension~$2$ (\cref{thm:top-fin-two-isGenericallyRigidInj}). This
  property passes along isomorphisms
  (\cref{lem:isGenericallyRigidInj-iso}), so $G$ has it too.

  Let $|V| \ge 3$. The reverse decomposition
  (\cref{thm:isLaman-exists-typeI-or-typeII-reverse}) gives a vertex
  $v$ of $G$ of degree two or three. Suppose first that $v$ has degree
  two, with neighbours $a$ and $b$, and that $G - v$ is Laman. The
  graph $G - v$ has $|V| - 1$ vertices, so by induction it is
  injectively generically rigid. Hence so is the graph obtained from
  $G - v$ by the Type~I move at $a$ and $b$
  (\cref{thm:typeI-isGenericallyRigidInj-two}). This graph is
  isomorphic to $G$, by the map that sends its new vertex to $v$ and
  fixes every other vertex (\cref{sec:henneberg-decomp}). So $G$ is
  injectively generically rigid (\cref{lem:isGenericallyRigidInj-iso}).

  Otherwise $v$ has degree three, with neighbours $a$, $b$ and $c$,
  where $a$ and $b$ are not adjacent, and $G - v + ab$ is Laman. By
  induction, $G - v + ab$ is injectively generically rigid. Hence so is
  the graph obtained from it by the Type~II move on $ab$ and $c$
  (\cref{thm:typeII-isGenericallyRigidInj-two}). This move deletes the
  edge $ab$ again, and the graph it gives is isomorphic to $G$ in the
  same way. So $G$ is injectively generically rigid, as before.
\end{proof}

\subsection{The \texorpdfstring{$\Rightarrow$}{=>} direction}

We now prove that every graph that is generically rigid in the plane
contains a Laman spanning subgraph. This is the necessity half of
Laman's theorem~\cite{laman1970}. We find the edges of the subgraph
among the rows of the rigidity map. Let $p$ be a framework on $V$ in
$\R^{2}$. A set $I$ of edges is row-independent at $p$ when the rows of
$R(G, p)$ at the edges of $I$ are linearly independent
(\cref{def:edgeSet-rowIndependent}). We prove two facts about such
sets.

First, if no three of the points $p_v$ are collinear, then the edges of
every row-independent set form a $(2,3)$-sparse graph
(\cref{lem:isSparse-of-rowIndependent-two}). The necessity of this
count for minimal rigidity goes back to Maxwell~\cite{maxwell1864}.
Whiteley~\cite{whiteley1996} proves the count in every dimension as his
$d$-Counting Lemma~11.1.3, and Jord\'an~\cite{jordan2016} as his
Lemma~1.3.1.

Second, if $p$ is infinitesimally rigid for $G$, then $R(G, p)$ has
rank at least $2|V| - 3$. So some $2|V| - 3$ edges are row-independent
at $p$ (see the proof of \cref{lem:exists-rowIndependent-edge-basis}).
Infinitesimal rigidity is an open condition
(\cref{lem:isInfinitesimallyRigid-eventually}). So an infinitesimally
rigid framework can be moved slightly to one that is still
infinitesimally rigid and has no three points collinear
(\cref{lem:exists-affinelySpanning-of-eventually}). There the
$2|V| - 3$ row-independent edges form a $(2,3)$-sparse graph with
$2|V| - 3$ edges, which is a Laman graph
(\cref{thm:isGenericallyRigid-exists-isLaman-le}).

Jord\'an (p.~39) calls the row-independent edge sets of a framework the
independent sets of its rigidity matroid. For a generic framework in
the plane, Laman characterized them: an edge set is row-independent if
and only if it forms a $(2,3)$-sparse graph (Jord\'an,
Theorem~2.2.1). The first fact above is one half of this
characterization, and Laman's theorem needs only this half. We prove
the converse in \cref{sec:rigidity-matroid}, where the two halves
identify the rigidity matroid in the plane with the $(2,3)$-count
matroid.

\subsubsection{Row-independent edge sets}

Let $p$ be a framework on $V$ in $\R^{d}$. The row of $R(G, p)$ at an
edge $e = uv$ is the linear functional
$p' \mapsto R(G, p)(p')(e) = \langle p_u - p_v,\, p'_u - p'_v \rangle$
on $(\R^{d})^{V}$. In the standard bases it is the row of Jord\'an's
rigidity matrix indexed by $e$.

\begin{definition}[Row-independent edge set at a placement]
  \label{def:edgeSet-rowIndependent}
  \lean{SimpleGraph.EdgeSetRowIndependent}
  \leanok
  \uses{def:rigidityMap}
  Let $p$ be a framework on $V$ in $\R^{d}$, and let
  $I \subseteq E(G)$. We say that $I$ is \emph{row-independent at $p$}
  if the linear functionals $p' \mapsto R(G, p)(p')(e)$ on
  $(\R^{d})^{V}$, for $e \in I$, are linearly independent.
\end{definition}

\subsubsection{The rows of the rigidity map}

To find a large row-independent edge set, we use that the row rank of
a matrix equals its column rank. For this we regard the rows of
$R(G, p)$ as elements of the dual space $((\R^{d})^{V})^{*}$ of linear
functionals on $(\R^{d})^{V}$. The \emph{transpose} of $R(G, p)$ is
the linear map $R(G, p)^{\top}$ from the dual space of $\R^{E(G)}$ to
$((\R^{d})^{V})^{*}$ that sends a functional $\varphi$ to
$\varphi \circ R(G, p)$. In the standard bases its matrix is the
transpose of the rigidity matrix. When $V$ is finite, $R(G, p)^{\top}$
has the same rank as $R(G, p)$.

\begin{definition}[Rigidity row at a placement]
  \label{def:rigidityRow}
  \lean{SimpleGraph.rigidityRow}
  \leanok
  \uses{def:rigidityMap}
  Let $p$ be a framework on $V$ in $\R^{d}$, and let $e$ be an edge of
  $G$. The \emph{rigidity row} of $G$ at $p$ indexed by $e$ is the
  linear functional $r_e \in ((\R^{d})^{V})^{*}$ given by
  $r_e(p') = R(G, p)(p')(e)$.
\end{definition}

\begin{lemma}[The rigidity rows span the image of the transpose]
  \label{lem:span-range-rigidityRow}
  \lean{SimpleGraph.span_range_rigidityRow}
  \leanok
  \uses{def:rigidityRow, def:rigidityMap}
  Let $V$ be finite, and let $p$ be a framework on $V$ in $\R^{d}$.
  Then the rigidity rows $r_e$, for $e \in E(G)$, span the image of
  $R(G, p)^{\top}$.
\end{lemma}
\begin{proof}
  \leanok
  For an edge $e$, let $e^{*}$ be the coordinate functional of
  $\R^{E(G)}$ at $e$. Then
  $r_e = e^{*} \circ R(G, p) = R(G, p)^{\top}(e^{*})$. Since $E(G)$ is
  finite, the functionals $e^{*}$ form a basis of the dual space of
  $\R^{E(G)}$. So their images under $R(G, p)^{\top}$ span its image.
\end{proof}

\begin{lemma}[Row-independence as linear independence of the rigidity rows]
  \label{lem:edgeSetRowIndependent-iff-linearIndepOn-rigidityRow}
  \lean{SimpleGraph.edgeSetRowIndependent_iff_linearIndepOn_rigidityRow}
  \leanok
  \uses{def:edgeSet-rowIndependent, def:rigidityRow}
  Let $p$ be a framework on $V$ in $\R^{d}$, and let
  $I \subseteq E(G)$. Then $I$ is row-independent at $p$ if and only if
  the rigidity rows $r_e$, for $e \in I$, are linearly independent in
  $((\R^{d})^{V})^{*}$.
\end{lemma}
\begin{proof}
  \leanok
  The rigidity row $r_e$ is the functional of
  \cref{def:edgeSet-rowIndependent} at the edge $e$. So both sides say
  that the same family of linear functionals is linearly independent.
\end{proof}
\begin{fmlnote}
  The formal definition of row-independence regards the rows as
  elements of the space of all real functions on $(\R^{d})^{V}$. The
  rigidity rows are elements of its linear subspace
  $((\R^{d})^{V})^{*}$, in which the transpose of $R(G, p)$ takes its
  values. The lemma passes between the two: a family of vectors of a
  subspace is linearly independent in the subspace exactly when it is
  linearly independent in the whole space.
\end{fmlnote}

\subsubsection{A lower bound on the rank}

The proof of the edge-count bound
(\cref{thm:isGenericallyRigid-card-mul-le}) first bounds the rank of
the rigidity map from below, and then bounds the rank by the number of
edges. We record the first step.

\begin{lemma}[Rank lower bound for generically rigid graphs]
  \label{lem:rigidityMap-finrank-range-ge-of-isGenericallyRigid}
  \lean{SimpleGraph.rigidityMap_finrank_range_ge_of_isGenericallyRigid}
  \leanok
  \uses{def:isGenericallyRigid, lem:framework-finrank,
    def:isInfinitesimallyRigid}
  Let $V$ be finite, and let $G$ be generically rigid in dimension $d$.
  Then there is a framework $p$ on $V$ in $\R^{d}$ with
  \[
    d\,|V| \;\le\; \operatorname{rank} R(G, p) + \frac{d(d+1)}{2}.
  \]
\end{lemma}
\begin{proof}
  \leanok
  Let $p$ be a framework that is infinitesimally rigid for $G$. By
  rank--nullity, $\dim (\R^{d})^{V}$ is the rank of $R(G, p)$ plus the
  dimension of its kernel. Here $\dim (\R^{d})^{V} = d\,|V|$
  (\cref{lem:framework-finrank}). The kernel has dimension at most
  $d(d+1)/2$, since $p$ is infinitesimally rigid
  (\cref{def:isInfinitesimallyRigid}).
\end{proof}

\begin{lemma}[A row-independent set of $2|V| - 3$ edges]
  \label{lem:exists-rowIndependent-edge-basis}
  \lean{SimpleGraph.exists_edgeSetRowIndependent_basis_dim_two}
  \leanok
  \uses{def:edgeSet-rowIndependent, def:rigidityRow,
    lem:rigidityMap-finrank-range-ge-of-isGenericallyRigid,
    lem:span-range-rigidityRow,
    lem:edgeSetRowIndependent-iff-linearIndepOn-rigidityRow}
  Let $V$ be finite, and let $G$ be generically rigid in
  dimension~$2$. Then there are a framework $p$ on $V$ in $\R^{2}$ and
  an edge set $I \subseteq E(G)$ with $|I| = \max(2|V| - 3,\, 0)$ that
  is row-independent at $p$.
\end{lemma}
\begin{proof}
  \leanok
  Let $p$ be a framework given by
  \cref{lem:rigidityMap-finrank-range-ge-of-isGenericallyRigid} at
  $d = 2$, so that $R(G, p)$ has rank at least $2|V| - 3$. The rigidity
  rows span the image of $R(G, p)^{\top}$
  (\cref{lem:span-range-rigidityRow}). This image has the same
  dimension as the image of $R(G, p)$, since the transpose has the same
  rank. So there is a set $B$ of edges whose rigidity rows are linearly
  independent and span the image of $R(G, p)^{\top}$, and
  $|B| \ge 2|V| - 3$.

  Let $I$ be a subset of $B$ with $\max(2|V| - 3,\, 0)$ elements. The
  rigidity rows $r_e$, for $e \in I$, are linearly independent. So $I$
  is row-independent at $p$
  (\cref{lem:edgeSetRowIndependent-iff-linearIndepOn-rigidityRow}).
\end{proof}

\subsubsection{An upper bound on the rank}

When the points $p_v$ affinely span $\R^{d}$, the kernel of $R(G, p)$
contains the space $T(p)$ of trivial motions, of dimension at least
$d(d+1)/2$ (\cref{lem:rigidityMap-ker-finrank-ge-of-affinelySpanning}).
By rank--nullity this bounds the rank from above. Jord\'an proves the
same bound for every framework on at least $d + 2$ vertices, as his
Lemma~1.2.1~\cite{jordan2016}. In the plane the bound gives the
sparsity of row-independent edge sets
(\cref{lem:isSparse-of-rowIndependent-two}).

\begin{lemma}[Rank upper bound for affinely spanning frameworks]
  \label{lem:rigidityMap-finrank-range-le-of-affinelySpanning}
  \lean{SimpleGraph.rigidityMap_finrank_range_le_of_affinelySpanning}
  \leanok
  \uses{lem:rigidityMap-ker-finrank-ge-of-affinelySpanning,
    lem:framework-finrank}
  Let $V$ be finite, and let $p$ be a framework on $V$ in $\R^{d}$
  whose points affinely span $\R^{d}$. Then
  \[
    \operatorname{rank} R(G, p) + \frac{d(d+1)}{2} \;\le\; d\,|V|.
  \]
\end{lemma}
\begin{proof}
  \leanok
  By rank--nullity, the rank of $R(G, p)$ plus the dimension of its
  kernel is $\dim (\R^{d})^{V}$, which is $d\,|V|$
  (\cref{lem:framework-finrank}). The kernel has dimension at least
  $d(d+1)/2$ (\cref{lem:rigidityMap-ker-finrank-ge-of-affinelySpanning}).
\end{proof}

\subsubsection{Sparsity from row-independence (Lov\'asz--Yemini, easy direction)}

The sparsity claim transports the rank upper bound from the induced
subframework $H[S]$ at $p|_{S}$ back to the row count of the
$H[S]$-rows in $\mathrm{RigidityMap}\,G\,p$. The two sides match
exactly because edges in $H[S]$ have both endpoints in $S$, so the
corresponding rows are supported on the columns indexed by $S$.
We need the witness placement to affinely span on every
size-$\ge 3$ subset; \cref{lem:exists-affinelySpanning-of-eventually}
below provides one at any dimension $d$ (used here at $d = 2$).

\begin{lemma}[Affinely-spanning perturbation of an eventually-true property, dimension $d$]
  \label{lem:exists-affinelySpanning-of-eventually}
  \lean{SimpleGraph.exists_affinelySpanning_of_eventually}
  \leanok
  Let $V$ be finite and $d \in \N$. Suppose $P$ is a property of
  placements $V \to \R^{d}$ that holds on a neighbourhood of some
  placement $p_{0}$. Then there exists a placement $p : V \to \R^{d}$
  satisfying $P$ such that the restriction $p|_{S}$ affinely spans
  $\R^{d}$ for every $S \subseteq V$ with $|S| \ge d + 1$.
\end{lemma}
\begin{proof}
  \leanok
  Choose an injection $\phi : V \to \R$ via the finite-type bijection
  $V \simeq \mathrm{Fin}\,|V|$ composed with the natural injection
  $\mathrm{Fin}\,n \hookrightarrow \R$, and set the moment-curve
  direction $w(v) := (\phi(v), \phi(v)^{2}, \ldots, \phi(v)^{d})
  \in \R^{d}$. For $t \in \R$ define the perturbed placement
  $p_{t} := p_{0} + t \cdot w$.

  The neighbourhood hypothesis on $P$, pulled back along the
  continuous map $t \mapsto p_{t}$, supplies an $\varepsilon > 0$
  such that $P(p_{t})$ holds for every $|t| < \varepsilon$.

  For each ordered $(d+1)$-tuple $(v_{0}, \ldots, v_{d})$ of distinct
  vertices, the affine-dependence determinant of the perturbed
  difference matrix
  \[
    \det\bigl(p_{t}(v_{1}) - p_{t}(v_{0}),\, \ldots,\,
      p_{t}(v_{d}) - p_{t}(v_{0})\bigr)
  \]
  is a polynomial in $t$ of degree at most $d$
  (\texttt{Polynomial.natDegree\_det\_X\_add\_C\_le}) whose top
  coefficient (\texttt{Polynomial.coeff\_det\_X\_add\_C\_card}) is the
  Vandermonde-difference determinant
  $\prod_{0 \le i < j \le d} (\phi(v_{j}) - \phi(v_{i}))$
  (\texttt{Matrix.det\_powerDifferences}), nonzero by injectivity of
  $\phi$. The per-tuple bad-$t$ set is therefore finite
  (\texttt{Polynomial.finite\_setOf\_isRoot}). The global bad set is
  the finite union of finitely many finite sets, hence finite, so any
  point of the (uncountable) interval $(0, \varepsilon)$ outside the
  bad set is a valid choice of $t$. The resulting $p_{t}$ satisfies
  $P$, and for every $S$ with $|S| \ge d + 1$ any $d + 1$ distinct
  elements of $S$ give an affinely-independent tuple under $p_{t}$,
  so $p_{t}|_{S}$ affinely spans $\R^{d}$.
\end{proof}
\begin{fmlnote}
  The moment-curve perturbation preserves any open property of the
  starting placement, so the lemma is stated for an arbitrary open
  property $P$ rather than for a fixed one. It is specialized twice:
  here to infinitesimal rigidity (with
  \cref{lem:isInfinitesimallyRigid-eventually} as the openness
  witness), and, for the $(\Rightarrow)$ direction of the matroid
  identification (\cref{thm:edgeSet-rowIndependent-iff-isSparse}), to
  row-independence (with \cref{lem:edgeSet-rowIndependent-eventually}).
  The single polymorphic statement avoids proving the moment-curve /
  Vandermonde-determinant construction twice.
\end{fmlnote}

\begin{lemma}[$(2, 3)$-sparsity from row-independence, dimension $2$]
  \label{lem:isSparse-of-rowIndependent-two}
  \lean{SimpleGraph.isSparse_of_edgeSetRowIndependent_dim_two}
  \leanok
  \uses{def:isSparse, def:edgeSet-rowIndependent,
    lem:rigidityMap-finrank-range-le-of-affinelySpanning,
    lem:exists-affinelySpanning-of-eventually}
  Let $V$ be finite, let $G$ be a graph on $V$, and let $p$ be a
  framework on $V$ in dimension $2$ such that $p|_{S}$ affinely spans
  for every $S \subseteq V$ with $|S| \ge 3$. If $I \subseteq E(G)$ is
  row-independent at $p$, then the spanning subgraph $H \le G$ with
  edge set $I$ is $(2, 3)$-sparse.
\end{lemma}
\begin{proof}
  \leanok
  Fix $S \subseteq V$. The edges of $H$ with both endpoints in $S$
  correspond bijectively to the edges of the induced subgraph $H[S]$;
  we show $|\edgesIn[H]{S}| + 3 \le 2 |S|$.

  For $|S| \le 1$ the proviso $3 \le 2|S|$ in $(2, 3)$-sparsity
  (\cref{def:isSparse}) is false, so the bound holds vacuously. For
  $|S| = 2$ the induced subgraph $H[S]$ has at most one edge (a simple
  graph on two vertices has at most one edge,
  \texttt{card\_edgeFinset\_le\_card\_choose\_two}), and
  $1 + 3 \le 4 = 2 |S|$.

  For $|S| \ge 3$: the rows of $\mathrm{RigidityMap}\,G\,p$ indexed
  by $\edgesIn[H]{S}$ are supported on the columns indexed by $S$
  (each row's nonzero entries are at the two endpoints of the
  corresponding edge, both of which lie in $S$). Hence row-independence
  of these rows in the full matrix transports through the
  framework-restriction map $\mathrm{Framework}\,V\,2 \to
  \mathrm{Framework}\,S\,2$ to row-independence of the
  $S$-indexed-column matrix, which is exactly the rigidity
  matrix $\mathrm{RigidityMap}\,(H[S])\,(p|_{S})$. The transport uses
  \texttt{LinearIndependent.of\_comp} on the dual-map side. By
  \cref{lem:rigidityMap-finrank-range-le-of-affinelySpanning} at
  $d = 2$ applied to $H[S]$ at the affinely-spanning placement
  $p|_{S}$, the rank of that matrix is at most $2 |S| - 3$. Since
  row-independence forces the row count $|\edgesIn[H]{S}|$ to equal
  the rank, $|\edgesIn[H]{S}| \le 2 |S| - 3$.
\end{proof}

\subsubsection{Assembly}

\begin{theorem}[Every generically rigid graph contains a Laman subgraph]
  \label{thm:isGenericallyRigid-exists-isLaman-le}
  \lean{SimpleGraph.IsGenericallyRigid.exists_isLaman_le}
  \leanok
  \uses{def:isGenericallyRigid, def:isLaman, def:isTight,
    lem:rigidityMap-finrank-range-ge-of-isGenericallyRigid,
    lem:exists-rowIndependent-edge-basis,
    lem:isSparse-of-rowIndependent-two,
    lem:exists-affinelySpanning-of-eventually,
    lem:isInfinitesimallyRigid-eventually}
  Let $V$ be finite with $|V| \ge 2$. If $G$ is generically rigid in
  dimension $2$, then there exists a Laman spanning subgraph
  $H \le G$.
\end{theorem}
\begin{proof}
  \leanok
  Choose a framework $p$ that is both infinitesimally rigid for $G$
  and affinely spanning on every subset of size $\ge 3$: apply
  \cref{lem:exists-affinelySpanning-of-eventually} to the open property
  of infinitesimal rigidity (\cref{lem:isInfinitesimallyRigid-eventually}),
  starting from an infinitesimally rigid placement supplied by the
  generic rigidity of $G$. At this $p$ the rank lower bound
  \cref{lem:rigidityMap-finrank-range-ge-of-isGenericallyRigid} gives
  $\dim\,\mathrm{range}\,(\mathrm{RigidityMap}\,G\,p) \ge 2 |V| - 3$, so
  as in \cref{lem:exists-rowIndependent-edge-basis} there is an edge set
  $I \subseteq E(G)$ with $|I| = 2 |V| - 3$ that is row-independent at
  $p$. Let $H \le G$ be the spanning subgraph with edge set $I$. By
  \cref{lem:isSparse-of-rowIndependent-two}, $H$ is $(2, 3)$-sparse;
  together with $|E(H)| = 2 |V| - 3$ this makes $H$ $(2, 3)$-tight,
  i.e.\ Laman.
\end{proof}

\subsection{Laman's theorem}

\begin{theorem}[Laman's theorem]
  \label{thm:isGenericallyRigid-two-iff-exists-isLaman-le}
  \lean{SimpleGraph.isGenericallyRigid_two_iff_exists_isLaman_le}
  \leanok
  \uses{thm:isLaman-isGenericallyRigid-two,
    thm:isGenericallyRigid-exists-isLaman-le,
    lem:isGenericallyRigid-mono}
  Let $V$ be finite with $|V| \ge 2$ and let $G$ be a graph on $V$.
  Then $G$ is generically rigid in dimension $2$ if and only if it
  contains a Laman spanning subgraph, i.e.\ there is a spanning
  subgraph $H \le G$ that is Laman.
\end{theorem}
\begin{proof}
  \leanok
  The $\Rightarrow$ direction is
  \cref{thm:isGenericallyRigid-exists-isLaman-le}. For $\Leftarrow$:
  given $H \le G$ with $H$ Laman,
  \cref{thm:isLaman-isGenericallyRigid-two} shows $H$ is generically
  rigid in dimension $2$, and monotonicity
  (\cref{lem:isGenericallyRigid-mono}) along $H \le G$ extends
  rigidity from $H$ to $G$.
\end{proof}
